\documentclass[12pt]{article}
\usepackage[utf8]{inputenc}
\usepackage{amsmath}
\usepackage{amssymb}
\usepackage{graphicx}
\usepackage{booktabs}
\usepackage{listings}
\usepackage{xcolor}
\usepackage{hyperref}
\usepackage{geometry}
\geometry{margin=1in}

\lstset{
    language=Python,
    basicstyle=\ttfamily\small,
    keywordstyle=\color{blue}\bfseries,
    commentstyle=\color{green!60!black},
    stringstyle=\color{red},
    numbers=left,
    numberstyle=\tiny\color{gray},
    stepnumber=1,
    numbersep=5pt,
    backgroundcolor=\color{gray!10},
    frame=single,
    breaklines=true,
    breakatwhitespace=false,
    tabsize=4,
    showstringspaces=false
}

\title{CSV Output Manager Documentation}
\author{}
\date{}

\begin{document}

\maketitle

\section{Introduction}

The \texttt{csv\_output\_manager.py} script provides the \texttt{CTCCOutputManager} class, which manages all CSV outputs for CTCC batch analysis, sensitivity studies, and Monte Carlo simulations. It coordinates data from all calculation modules into a unified \texttt{batch\_summary.csv} and individual module-specific CSVs.

\section{Method Overview}

\subsection{Purpose}

The CSV output manager:
\begin{enumerate}
    \item Maintains a single \texttt{batch\_summary.csv} with one row per scenario
    \item Creates module-specific CSVs with detail and summary rows
    \item Automatically prepends technical parameters to all CSV rows
    \item Manages scenario IDs for coordinated batch runs
    \item Updates existing rows instead of creating duplicates
    \item Handles both detail rows (terrain-specific, component-specific) and summary rows
\end{enumerate}

\subsection{Architecture}

The manager uses a two-tier CSV structure:
\begin{itemize}
    \item \textbf{Batch Summary}: One comprehensive row per scenario with all metrics
    \item \textbf{Module CSVs}: Detail rows (e.g., by terrain, by component) plus summary row
\end{itemize}

All CSVs automatically include technical parameters as the first columns for easy filtering and comparison.

\subsection{Integration with Other Scripts}

This script is used by:
\begin{itemize}
    \item All cost/benefit calculation scripts --- via \texttt{smart\_output.CTCCOutputManager} (dispatches to \texttt{CSVOutputManager} or \texttt{JSONOutputManager} by mode); each calls \texttt{add\_*} methods and \texttt{write\_batch\_summary}
    \item \texttt{bcr\_calculator.py} --- Reads \texttt{batch\_summary.csv} and writes BCR columns
    \item \texttt{ctcc.py} --- Sets \texttt{CTCC\_SCENARIO\_ID} and coordinates final summary
\end{itemize}

\textbf{Dependencies:}
\begin{itemize}
    \item \texttt{path\_config.py} --- \texttt{OUTPUTS\_DIR}, \texttt{YAMLS\_DIR}
    \item \texttt{smart\_loaders.py} --- \texttt{load\_project\_technical\_details}, \texttt{load\_physical\_details}, \texttt{load\_financing\_social\_discount\_rate}, \texttt{get\_project\_data\_raw} (for \texttt{load\_technical\_details}; supports YAML and JSON input mode)
\end{itemize}

\section{Key Functions}

\subsection{\texttt{\_\_init\_\_(output\_dir, scenario\_id)}}

Initializes the output manager.

\textbf{Parameters:}
\begin{itemize}
    \item \texttt{output\_dir} (str): Base directory for outputs (default: \texttt{"../outputs"})
    \item \texttt{scenario\_id} (str): Unique scenario identifier (optional)
\end{itemize}

\textbf{Key Logic:}
\begin{lstlisting}
# Scenario ID: explicit > CTCC_SCENARIO_ID env > auto-generated (with microseconds)
self.scenario_id = (scenario_id or os.environ.get("CTCC_SCENARIO_ID")
    or datetime.now().strftime("%Y%m%d_%H%M%S_%f"))
self.timestamp = datetime.now().isoformat()
technical_details = self.load_technical_details()
self.batch_summary_data = {**technical_details, "scenario_id": self.scenario_id, "timestamp": self.timestamp}
\end{lstlisting}

\subsection{\texttt{load\_technical\_details()}}

Loads technical parameters using centralized loaders (\texttt{smart\_loaders}: \texttt{load\_project\_technical\_details}, \texttt{load\_physical\_details}, \texttt{load\_financing\_social\_discount\_rate}, \texttt{get\_project\_data\_raw}) to include in CSV outputs. Supports both YAML and JSON input modes.

\textbf{Returns:}
\begin{itemize}
    \item Dictionary of technical parameters (project name, construction type, capacity, line length, discount rate, timeline, etc.)
\end{itemize}

\textbf{Technical Parameters Loaded:}
\begin{itemize}
    \item Project identification (name)
    \item Core specs (construction\_type, ac\_dc, capacity\_mw, conductor\_type, line\_length\_miles, line\_utilization)
    \item Converter details (converter\_type, number\_of\_converters)
    \item Reconductoring details (reconductoring, old\_capacity\_mw, old\_conductor\_type, old\_ac\_dc)
    \item Financial parameters (baseline\_electricity\_price\_per\_mwh, social\_discount\_rate)
    \item Timeline (construction\_years, delay\_years, project\_lifetime\_years)
\end{itemize}

\subsection{\texttt{write\_batch\_summary()}}

Writes or updates the batch summary to CSV. If a row with the same \texttt{scenario\_id} exists, it updates that row; otherwise, it appends a new row.

\textbf{Key Logic:}
\begin{lstlisting}
# Check if scenario_id already exists
scenario_row_idx = None
for idx, row in enumerate(existing_rows):
    if row.get('scenario_id') == self.scenario_id:
        scenario_row_idx = idx
        break

# Update existing row or append new one
if scenario_row_idx is not None:
    existing_rows[scenario_row_idx].update(self.batch_summary_data)
else:
    existing_rows.append(self.batch_summary_data)
\end{lstlisting}

\subsection{\texttt{write\_module\_csv(module\_name, detail\_rows, summary\_row, append)}}

Writes module-specific CSV with optional detail and summary rows. Automatically prepends technical parameters to each row.

\textbf{Parameters:}
\begin{itemize}
    \item \texttt{module\_name} (str): Name of the module (e.g., \texttt{'wildfire\_costs'})
    \item \texttt{detail\_rows} (list): List of dicts for detail rows (optional)
    \item \texttt{summary\_row} (dict): Dict for summary row (optional)
    \item \texttt{append} (bool): If True, append to existing file; if False, overwrite
\end{itemize}

\textbf{Key Logic:}
\begin{lstlisting}
# Extract technical parameters
tech_params = {k: self.batch_summary_data.get(k, '') 
               for k in tech_param_keys}
tech_params['scenario_id'] = self.scenario_id
tech_params['timestamp'] = self.timestamp

# Prepend technical parameters to each row
for row in all_rows:
    enriched_row = {**tech_params, **row}  # Tech params first
    enriched_rows.append(enriched_row)
\end{lstlisting}

\subsection{Module-Specific Add Methods}

Each calculation module has a corresponding \texttt{add\_*} method:

\begin{itemize}
    \item \texttt{add\_build\_costs(results)} - Build costs
    \item \texttt{add\_row\_costs(results)} - Right-of-way costs
    \item \texttt{add\_environmental\_mitigation(results)} - Environmental mitigation
    \item \texttt{add\_delay\_costs(results)} - Delay costs
    \item \texttt{add\_insurance\_costs(results)} - Insurance costs
    \item \texttt{add\_wildfire\_costs(results)} - Wildfire costs (with detail rows by terrain)
    \item \texttt{add\_outage\_costs(results)} - Outage costs (with detail rows by terrain)
    \item \texttt{add\_oandm\_costs(results)} - O\&M costs (with detail rows by component)
    \item \texttt{add\_emissions\_costs(results)} - Emissions costs (with detail rows by pollutant)
    \item \texttt{add\_line\_loss\_costs(results)} - Line loss costs/benefits
    \item \texttt{add\_congestion\_curtailment(results)} - Congestion and curtailment benefits/costs
    \item \texttt{add\_bcr\_metrics(bcr\_results)} - BCR metrics
\end{itemize}

Each method:
\begin{enumerate}
    \item Updates \texttt{batch\_summary\_data} with summary metrics
    \item Creates detail rows (if applicable)
    \item Creates summary row
    \item Calls \texttt{write\_module\_csv()} to write the module CSV
\end{enumerate}

\subsection{\texttt{add\_congestion\_curtailment(results)}}

Special method that distinguishes benefits from costs in detail rows.

\textbf{Detail Rows Created:}
\begin{itemize}
    \item \textbf{Benefits}: Congestion reduction (full and haircut), Curtailment reduction (full and haircut)
    \item \textbf{Costs}: Congestion during delay, Curtailment during delay, Residual congestion
\end{itemize}

Each detail row includes:
\begin{itemize}
    \item \texttt{benefit\_or\_cost}: "benefit" or "cost"
    \item \texttt{constraint\_type}: "congestion", "curtailment", "congestion\_delay", etc.
    \item \texttt{value\_type}: "full", "haircut", or "NA"
    \item Annual, nominal, and PV values
\end{itemize}

\subsection{\texttt{add\_bcr\_metrics(bcr\_results)}}

Adds BCR metrics to the batch summary.

\textbf{Expected Keys in bcr\_results:}
\begin{itemize}
    \item \texttt{bcr\_system}, \texttt{bcr\_capital}, \texttt{bcr\_haircut}, \texttt{bcr\_excluding\_risk}
    \item \texttt{net\_benefit\_pv}, \texttt{net\_benefit\_nominal}, \texttt{net\_benefit\_excluding\_risk\_pv}
    \item \texttt{total\_benefits\_pv}, \texttt{total\_costs\_pv}, etc.
\end{itemize}

\section{Example}

The following example demonstrates typical usage:

\begin{lstlisting}
# When called from build_costs.py:
csv_manager = CTCCOutputManager()

# Add build cost results
results = {
    "total_nominal": 100000000,
    "total_afudc": 105000000,
    "total_pv": 95000000,
    "conductor_nominal": 60000000,
    "structure_nominal": 30000000,
    "converter_nominal": 10000000,
}
csv_manager.add_build_costs(results)
csv_manager.write_batch_summary()

# This:
# 1. Updates batch_summary_data with build cost metrics
# 2. Creates build_costs.csv with:
#    - Technical parameters (first 20 columns)
#    - Summary row with all build cost metrics
# 3. Updates batch_summary.csv with build cost columns
\end{lstlisting}

\textbf{Multi-Module Example:}

\begin{lstlisting}
# When called from ctcc.py (coordinated run):
import os
os.environ['CTCC_SCENARIO_ID'] = '20250101_120000'

# Each script creates its own manager instance
# All use the same scenario_id from environment variable

# build_costs.py:
csv_manager = CTCCOutputManager()  # Gets scenario_id from env
csv_manager.add_build_costs(results)
csv_manager.write_batch_summary()  # Creates/updates row

# row_costs.py:
csv_manager = CTCCOutputManager()  # Same scenario_id
csv_manager.add_row_costs(results)
csv_manager.write_batch_summary()  # Updates same row

# Result: batch_summary.csv has one row with all metrics
\end{lstlisting}

\textbf{Detail Rows Example:}

\begin{lstlisting}
# When called from wildfire_costs.py:
csv_manager = CTCCOutputManager()
results = {
    "EAL": 500000,
    "lambda_total": 0.01,
    "lambda_by_terrain": {
        "forested": {"miles": 40, "events_per_year": 0.006, ...},
        "farmland": {"miles": 60, "events_per_year": 0.004, ...},
    },
    ...
}
csv_manager.add_wildfire_costs(results)

# This creates wildfire_costs.csv with:
# - Technical parameters (first 20 columns)
# - Detail rows (one per terrain) with terrain-specific metrics
# - Summary row with aggregated metrics
\end{lstlisting}

\section{Key Implementation Details}

\subsection{Scenario ID Management}

The manager uses a three-tier priority for scenario IDs:
\begin{enumerate}
    \item Explicit parameter (if provided)
    \item Environment variable \texttt{CTCC\_SCENARIO\_ID} (for coordinated batch runs)
    \item Auto-generated timestamp (if neither above)
\end{enumerate}

This ensures all scripts in a single CTCC run share the same scenario ID.

\subsection{Row Update vs. Append}

The \texttt{write\_batch\_summary()} method checks if a row with the same \texttt{scenario\_id} exists. If it does, it updates that row; otherwise, it appends a new row. This allows multiple scripts to contribute to the same scenario row.

\subsection{Technical Parameters Prepend}

All module CSVs automatically include technical parameters as the first columns. This is done in \texttt{write\_module\_csv()} by prepending the technical parameters dictionary to each row.

\subsection{Fieldname Management}

When updating \texttt{batch\_summary.csv}, the manager:
\begin{enumerate}
    \item Reads existing fieldnames
    \item Merges with new fieldnames (existing first, then new)
    \item Writes all rows with the merged fieldnames
\end{enumerate}

This ensures backward compatibility when new columns are added.

\subsection{Detail vs. Summary Rows}

Module CSVs can include:
\begin{itemize}
    \item \textbf{Detail rows}: Granular data (e.g., by terrain, by component, by pollutant)
    \item \textbf{Summary row}: Aggregated totals
\end{itemize}

Each row is marked with \texttt{row\_type}: "detail" or "summary".

\subsection{Path Resolution}

The manager uses dynamic path resolution to find YAML files:
\begin{lstlisting}
script_dir = os.path.dirname(os.path.abspath(__file__))
yaml_dir = os.path.join(os.path.dirname(script_dir), "yamls")
\end{lstlisting}

This ensures the script works regardless of the current working directory.

\subsection{Error Handling}

The \texttt{load\_technical\_details()} method includes try/except to handle missing YAML files gracefully, returning an empty dictionary and printing a warning.

\end{document}

